% ============================================================
%  REPORTE — Reconocimiento de Patrones
%  Universidad Tecnológica de la Mixteca (UTM)
% ============================================================
\documentclass[12pt, letterpaper]{article}

% ----- Codificación y lenguaje -----
\usepackage[utf8]{inputenc}
\usepackage[T1]{fontenc}
\usepackage[spanish]{babel}

% ----- Márgenes -----
\usepackage[top=2.5cm, bottom=2.5cm, left=3cm, right=2.5cm]{geometry}

% ----- Fuente profesional -----
\usepackage{lmodern}

% ----- Tablas y gráficos -----
\usepackage{booktabs}
\usepackage{array}
\usepackage{multirow}
\usepackage{graphicx}
\usepackage{float}
\usepackage{xcolor}
\usepackage{colortbl}

% ----- Matemáticas -----
\usepackage{amsmath}

% ----- Interlineado -----
\usepackage{setspace}
\onehalfspacing

% ----- Cabecera / pie -----
\usepackage{fancyhdr}
\pagestyle{fancy}
\fancyhf{}
\fancyhead[L]{\small Reconocimiento de Patrones}
\fancyhead[R]{\small UTM — Instituto de Computación}
\fancyfoot[C]{\thepage}
\renewcommand{\headrulewidth}{0.4pt}

% ----- Hipervínculos -----
\usepackage[hidelinks]{hyperref}

% ----- Color para filas alternadas en tablas -----
\definecolor{rowgray}{gray}{0.93}

% ============================================================
\begin{document}

% ============================================================
%  PORTADA
% ============================================================
\begin{titlepage}
    \centering
    \vspace*{1.5cm}

    {\Large \textbf{Universidad Tecnológica de la Mixteca}}\\[0.3cm]
    {\large Instituto de Computación}\\[0.2cm]
    {\large Ingeniería en Computación}

    \vspace{1.5cm}
    \rule{\linewidth}{0.5pt}\\[0.4cm]
    {\LARGE \textbf{Clasificación de Niveles de Ansiedad y Estrés\\
    mediante Naive Bayes y Árboles de Decisión}}\\[0.4cm]
    \rule{\linewidth}{0.5pt}

    \vspace{1.5cm}
    {\large \textbf{Materia:} Reconocimiento de Patrones}\\[0.6cm]

    \vspace{1cm}
    {\large \textbf{Estudiante:}}\\[0.3cm]
    {\large José Ramón Aragón Toledo}

    \vfill
    {\large Huajuapan de León, Oaxaca \hfill Abril, 2026}
\end{titlepage}

% ============================================================
%  ÍNDICE
% ============================================================
\tableofcontents
\newpage

% ============================================================
\section{Introducción}
% ============================================================

El presente reporte documenta el proceso de clasificación de niveles de estrés
utilizando el conjunto de datos \textit{StressLevelDataset}, el cual contiene
1\,100 registros de estudiantes con 20 variables fisiológicas, psicológicas y
ambientales, y una variable objetivo \texttt{stress\_level} con tres clases:
\textbf{Bajo (0)}, \textbf{Medio (1)} y \textbf{Alto (2)}.

El objetivo es comparar el desempeño de dos algoritmos de clasificación —
\textit{Gaussian Naive Bayes} y \textit{Decision Tree Classifier} — bajo dos
esquemas de características:

\begin{enumerate}
    \item \textbf{Tarea 1:} Clasificación utilizando las 20 variables disponibles.
    \item \textbf{Tarea 2:} Clasificación utilizando únicamente las 10 mejores
          variables según el criterio de \textit{Información Mutua} (MMI).
\end{enumerate}

Todos los experimentos se ejecutaron con \texttt{Python 3.12} y la biblioteca
\texttt{scikit-learn 1.8}.

% ============================================================
\section{Preprocesamiento de Datos}
% ============================================================

\subsection{Carga y verificación de calidad}

El dataset fue cargado con \texttt{pandas}. La inspección formal arrojó los
siguientes resultados:

\begin{itemize}
    \item \textbf{Valores nulos:} 0 en las 21 columnas.
    \item \textbf{Filas duplicadas:} 0.
    \item \textbf{Tipos de datos:} todos enteros (\texttt{int64}), sin necesidad
          de codificación categórica.
\end{itemize}

\subsection{Distribución de la variable objetivo}

Las tres clases presentan una distribución prácticamente uniforme, lo que
garantiza que los resultados no estén sesgados por desbalance de clases
(ver Tabla~\ref{tab:distribucion}).

\begin{table}[H]
    \centering
    \caption{Distribución de la variable objetivo \texttt{stress\_level}}
    \label{tab:distribucion}
    \begin{tabular}{lcc}
        \toprule
        \textbf{Clase} & \textbf{Frecuencia} & \textbf{Porcentaje (\%)} \\
        \midrule
        \rowcolor{rowgray} Bajo (0)  & 373 & 33.91 \\
        Medio (1)                    & 358 & 32.55 \\
        \rowcolor{rowgray} Alto (2)  & 369 & 33.55 \\
        \midrule
        \textbf{Total} & 1\,100 & 100.00 \\
        \bottomrule
    \end{tabular}
\end{table}

\subsection{Escalado y división de datos}

Para normalizar las escalas de las características se aplicó
\textbf{StandardScaler} de \texttt{scikit-learn}. El escalado se ajustó
\emph{únicamente} sobre el conjunto de entrenamiento y posteriormente se
transformaron ambos conjuntos, evitando así la fuga de información
(\textit{data leakage}).

La división \textbf{Hold-out} se realizó con los siguientes parámetros:

\begin{itemize}
    \item Proporción: \textbf{70\% entrenamiento / 30\% prueba}.
    \item \texttt{random\_state = 42} para garantizar replicabilidad.
    \item \texttt{stratify = y} para mantener la proporción de clases en ambos
          subconjuntos.
\end{itemize}

\begin{table}[H]
    \centering
    \caption{Tamaños de los conjuntos resultantes}
    \label{tab:split}
    \begin{tabular}{lcc}
        \toprule
        \textbf{Conjunto} & \textbf{Muestras} & \textbf{Variables} \\
        \midrule
        \rowcolor{rowgray} Entrenamiento (\textit{X\_train}) & 770 & 20 \\
        Prueba (\textit{X\_test})                            & 330 & 20 \\
        \bottomrule
    \end{tabular}
\end{table}

% ============================================================
\section{Tarea 1: Clasificación con Todas las Variables}
% ============================================================

\subsection{Modelos utilizados}

\begin{itemize}
    \item \textbf{Gaussian Naive Bayes} (\texttt{GaussianNB}): asume
          independencia condicional entre características y distribución
          gaussiana por clase.
    \item \textbf{Decision Tree Classifier} con \texttt{max\_depth = 6}:
          límite de profundidad para controlar el sobreajuste
          (\textit{overfitting}).
\end{itemize}

\subsection{Resultados}

La Tabla~\ref{tab:tarea1} muestra las métricas globales (ponderadas por clase)
obtenidas en el conjunto de prueba.

\begin{table}[H]
    \centering
    \caption{Métricas — Tarea 1 (20 variables)}
    \label{tab:tarea1}
    \begin{tabular}{lcccc}
        \toprule
        \textbf{Modelo} & \textbf{Accuracy} & \textbf{Precision} &
        \textbf{Recall} & \textbf{F1-Score} \\
        \midrule
        \rowcolor{rowgray}
        Gaussian Naive Bayes       & 88.18\% & 0.8990 & 0.8818 & 0.8837 \\
        Decision Tree (depth = 6)  & 88.48\% & 0.8853 & 0.8848 & 0.8849 \\
        \bottomrule
    \end{tabular}
\end{table}

\subsection{Matrices de confusión}

\begin{table}[H]
    \centering
    \caption{Matriz de Confusión — Gaussian Naive Bayes (Tarea 1)}
    \label{tab:cm_gnb_t1}
    \begin{tabular}{lccc}
        \toprule
        \textbf{Real $\backslash$ Predicho} & \textbf{Bajo (0)} &
        \textbf{Medio (1)} & \textbf{Alto (2)} \\
        \midrule
        \rowcolor{rowgray} Bajo (0)  &  94 &  0 & 18 \\
        Medio (1)                    &   3 & 90 & 14 \\
        \rowcolor{rowgray} Alto (2)  &   4 &  0 & 107 \\
        \bottomrule
    \end{tabular}
\end{table}

\begin{table}[H]
    \centering
    \caption{Matriz de Confusión — Decision Tree (Tarea 1)}
    \label{tab:cm_dt_t1}
    \begin{tabular}{lccc}
        \toprule
        \textbf{Real $\backslash$ Predicho} & \textbf{Bajo (0)} &
        \textbf{Medio (1)} & \textbf{Alto (2)} \\
        \midrule
        \rowcolor{rowgray} Bajo (0)  & 98 &  5 &  9 \\
        Medio (1)                    &  5 & 95 &  7 \\
        \rowcolor{rowgray} Alto (2)  &  6 &  6 & 99 \\
        \bottomrule
    \end{tabular}
\end{table}

% ============================================================
\section{Tarea 2: Clasificación con Ranking por Información Mutua}
% ============================================================

\subsection{Selección de características}

La \textbf{Información Mutua} (\textit{Mutual Information}, MMI) cuantifica la
dependencia estadística entre cada variable independiente y la variable
objetivo. Se calculó mediante \texttt{mutual\_info\_classif} de
\texttt{scikit-learn}, ajustándose únicamente sobre el conjunto de
entrenamiento para evitar sesgo.

El Top~5 del ranking se presenta en la Tabla~\ref{tab:ranking_mmi} y el
conjunto completo se encuentra en el archivo \texttt{results/ranking\_mmi.csv}.

\begin{table}[H]
    \centering
    \caption{Top 5 — Ranking por Información Mutua}
    \label{tab:ranking_mmi}
    \begin{tabular}{clc}
        \toprule
        \textbf{Posición} & \textbf{Característica} & \textbf{Score MMI} \\
        \midrule
        \rowcolor{rowgray} 1 & blood\_pressure               & 0.7576 \\
        2                    & future\_career\_concerns      & 0.6764 \\
        \rowcolor{rowgray} 3 & sleep\_quality                & 0.6722 \\
        4                    & bullying                      & 0.6627 \\
        \rowcolor{rowgray} 5 & self\_esteem                  & 0.6488 \\
        \bottomrule
    \end{tabular}
\end{table}

Para el entrenamiento se utilizaron las \textbf{Top 10} características
(hasta \texttt{teacher\_student\_relationship}, score = 0.5429),
reduciendo el espacio de características en un \textbf{50\%}.

\subsection{Resultados con Top-10 MMI}

\begin{table}[H]
    \centering
    \caption{Métricas — Tarea 2 (Top-10 MMI)}
    \label{tab:tarea2}
    \begin{tabular}{lcccc}
        \toprule
        \textbf{Modelo} & \textbf{Accuracy} & \textbf{Precision} &
        \textbf{Recall} & \textbf{F1-Score} \\
        \midrule
        \rowcolor{rowgray}
        Gaussian Naive Bayes       & 87.27\% & 0.8943 & 0.8727 & 0.8752 \\
        Decision Tree (depth = 6)  & 86.97\% & 0.8697 & 0.8697 & 0.8696 \\
        \bottomrule
    \end{tabular}
\end{table}

\subsection{Comparación Tarea 1 vs.\ Tarea 2}

\begin{table}[H]
    \centering
    \caption{Comparación de Accuracy y F1-Score entre ambas tareas}
    \label{tab:comparacion}
    \begin{tabular}{lcccc}
        \toprule
        \multirow{2}{*}{\textbf{Modelo}} &
        \multicolumn{2}{c}{\textbf{Accuracy}} &
        \multicolumn{2}{c}{\textbf{F1-Score (weighted)}} \\
        \cmidrule(lr){2-3} \cmidrule(lr){4-5}
        & \textbf{Tarea 1} & \textbf{Tarea 2} &
          \textbf{Tarea 1} & \textbf{Tarea 2} \\
        \midrule
        \rowcolor{rowgray}
        Gaussian Naive Bayes      & 88.18\% & 87.27\% & 0.8837 & 0.8752 \\
        Decision Tree (depth = 6) & 88.48\% & 86.97\% & 0.8849 & 0.8696 \\
        \bottomrule
    \end{tabular}
\end{table}

% ============================================================
\section{Análisis y Conclusiones}
% ============================================================

Los resultados obtenidos permiten establecer las siguientes conclusiones:

\begin{enumerate}

    \item \textbf{Eficiencia del ranking MMI.}
    La reducción del espacio de características de 20 a 10 variables
    (50\,\%) provocó una disminución del \textit{Accuracy} de apenas
    $\approx\!1\%$ en ambos modelos (de 88.18\,\% a 87.27\,\% en
    Naive Bayes; de 88.48\,\% a 86.97\,\% en el Árbol de Decisión).
    Este resultado demuestra que el ranking por Información Mutua identifica
    correctamente las variables más informativas y que las 10 variables
    descartadas aportaban escasa información discriminante.

    \item \textbf{Recall perfecto de Naive Bayes para la clase ``Alto''.}
    En ambas tareas, el modelo Gaussian Naive Bayes alcanzó un
    \textbf{Recall = 0.96} para la clase \textit{Alto (2)}, lo que significa
    que identificó correctamente el 96\,\% de los casos de estrés alto.
    Esta propiedad es especialmente valiosa en aplicaciones clínicas o
    educativas donde un falso negativo (no detectar estrés alto) tiene mayor
    costo que un falso positivo.

    \item \textbf{Balanceo del árbol de decisión.}
    El Decision Tree distribuyó los errores de forma más homogénea entre las
    tres clases, obteniendo métricas de Precision, Recall y F1 más uniformes
    (diferencia máxima de 0.04 entre clases en Tarea 1), mientras que
    Naive Bayes concentró sus errores en las clases ``Bajo'' y ``Medio''.

    \item \textbf{Rendimiento general.}
    Ambos modelos superaron el 86\,\% de \textit{Accuracy} incluso con la
    mitad de las variables, lo que confirma que el dataset presenta patrones
    bien definidos y que los algoritmos seleccionados son apropiados para
    este problema de clasificación multiclase.

\end{enumerate}

% ============================================================
%  BIBLIOGRAFÍA MÍNIMA
% ============================================================
\newpage
\begin{thebibliography}{9}

\bibitem{mitchell1997}
    Mitchell, T.~M. (1997).
    \textit{Machine Learning}.
    McGraw-Hill.

\bibitem{sklearn}
    Pedregosa, F. et al. (2011).
    Scikit-learn: Machine Learning in Python.
    \textit{Journal of Machine Learning Research}, 12, 2825--2830.

\bibitem{cover2006}
    Cover, T.~M. \& Thomas, J.~A. (2006).
    \textit{Elements of Information Theory} (2nd ed.).
    Wiley-Interscience.

\end{thebibliography}

\end{document}

% ============================================================
%  RESUMEN EN MARKDOWN (para PowerPoint / Excel)
%  Copia desde aquí hacia abajo
% ============================================================

% ## Tarea 1 — Todas las variables (20 features)
%
% | Modelo                   | Accuracy | Precision | Recall | F1-Score |
% |--------------------------|----------|-----------|--------|----------|
% | Gaussian Naive Bayes     | 88.18 %  | 0.8990    | 0.8818 | 0.8837   |
% | Decision Tree (depth=6)  | 88.48 %  | 0.8853    | 0.8848 | 0.8849   |
%
% ---
%
% ## Tarea 2 — Top-10 MMI (10 features)
%
% | Modelo                   | Accuracy | Precision | Recall | F1-Score |
% |--------------------------|----------|-----------|--------|----------|
% | Gaussian Naive Bayes     | 87.27 %  | 0.8943    | 0.8727 | 0.8752   |
% | Decision Tree (depth=6)  | 86.97 %  | 0.8697    | 0.8697 | 0.8696   |
%
% ---
%
% ## Ranking MMI — Top 5
%
% | # | Característica            | Score MMI |
% |---|---------------------------|-----------|
% | 1 | blood_pressure            | 0.7576    |
% | 2 | future_career_concerns    | 0.6764    |
% | 3 | sleep_quality             | 0.6722    |
% | 4 | bullying                  | 0.6627    |
% | 5 | self_esteem               | 0.6488    |
%
% ---
%
% ## Comparación Tarea 1 vs Tarea 2
%
% | Modelo                   | T1 Acc   | T2 Acc   | Δ Acc   | T1 F1  | T2 F1  | Δ F1    |
% |--------------------------|----------|----------|---------|--------|--------|---------|
% | Gaussian Naive Bayes     | 88.18 %  | 87.27 %  | −0.91 % | 0.8837 | 0.8752 | −0.0085 |
% | Decision Tree (depth=6)  | 88.48 %  | 86.97 %  | −1.51 % | 0.8849 | 0.8696 | −0.0153 |
